\unitopen{17}
\begin{candos}
\item użyć \de{Präteritum} czasowników \de{sein}, \de{haben} i modalnych (\de{war, hatte, konnte, musste});
\item nazwać części ciała i opisać dolegliwości u lekarza;
\item umówić się na wizytę i zrozumieć proste zalecenia;
\item rozumieć, dlaczego Perfekt dominuje w mowie, a Präteritum w kilku słowach.
\end{candos}

\section{Słownictwo}
\vocabtable{17}

\section{Präteritum czterech najważniejszych grup}
W mowie przeszłość to Perfekt, z \emph{trzema} wyjątkami, które mówisz w Präteritum, bo brzmią naturalnie:
\de{sein}, \de{haben} i czasowniki modalne.

\begin{center}\small
\begin{tabular}{@{}lllll@{}}
\toprule
 & \de{sein} & \de{haben} & \de{können} & \de{müssen}\\
\midrule
\de{ich} & \de{war} & \de{hatte} & \de{konnte} & \de{musste}\\
\de{du} & \de{warst} & \de{hattest} & \de{konntest} & \de{musstest}\\
\de{er/sie/es} & \de{war} & \de{hatte} & \de{konnte} & \de{musste}\\
\de{wir} & \de{waren} & \de{hatten} & \de{konnten} & \de{mussten}\\
\de{ihr} & \de{wart} & \de{hattet} & \de{konntet} & \de{musstet}\\
\de{sie/Sie} & \de{waren} & \de{hatten} & \de{konnten} & \de{mussten}\\
\bottomrule
\end{tabular}
\end{center}
Pozostałe modalne: \de{wollte}, \de{durfte}, \de{sollte} (bez przegłosu). Nawiązanie: \de{Ich war gestern krank.}
\de{Ich hatte Fieber.} \de{Ich konnte nicht arbeiten.} \de{Ich musste zum Arzt gehen.}

\section{U lekarza}
Skargi: \de{Mir tut der Kopf weh.} \de{Ich habe Kopfschmerzen / Bauchschmerzen / Rückenschmerzen.} \de{Ich habe Fieber und
Husten.} \de{Seit drei Tagen.} Lekarz: \de{Wo tut es weh?} \de{Seit wann haben Sie das?} \de{Nehmen Sie diese Tabletten
dreimal täglich.} \de{Sie müssen sich ausruhen.}

\begin{falle}
\de{Mir tut der Kopf weh} używa celownika (\de{mir}) i rodzajnika zamiast przymiotnika dzierżawczego (nie \de{mein Kopf}).
\de{Ich habe Schmerzen} (nie \de{Ich bin Schmerzen}). Drugi błąd: stosowanie Präteritum wszędzie (\de{Ich ging gestern
ins Kino}); w rozmowie powiedz \de{Ich bin gestern ins Kino gegangen}. Trzeci: \de{krank sein} znaczy „być chorym”,
a „zachorować” to \de{krank werden}.
\end{falle}

\begin{szwajcaria}
Ubezpieczenie zdrowotne jest w Szwajcarii \emph{obowiązkowe} i prywatne: wybierasz \de{Krankenkasse}
\verify{(w ciągu trzech miesięcy od zamieszkania; sprawdź na stronie władz federalnych i u kantonu)}. Płacisz
\de{Prämie}, \de{Franchise} (udział własny) i \de{Selbstbehalt}. Pierwszym punktem kontaktu jest zwykle
\de{Hausarzt}; pogotowie to \de{144}, a numer ogólny alarmowy \de{112}. Szpital to \de{das Spital}.
\end{szwajcaria}

\begin{praca}
Zwolnienie: \de{Ich bin heute krank und kann nicht arbeiten.} \de{Ich war gestern beim Arzt.} \de{Ich hatte Fieber.}
\de{Ich melde mich krank.} W CH pracodawca zwykle wymaga zaświadczenia lekarskiego od drugiego lub
trzeciego dnia \verify{(zależy od umowy i zakładu pracy)}.
\end{praca}

\section{Dialogi}
\textbf{Dialog 1 --- u lekarza}
\begin{dialog}
\turn{A}{Guten Tag. Was fehlt Ihnen?}{Dzień dobry. Co Panu dolega?}
\turn{B}{Ich habe seit drei Tagen Fieber und Husten. Mir tut der Hals weh.}{Od trzech dni mam gorączkę i kaszel. Boli mnie gardło.}
\turn{A}{Waren Sie schon bei einem Arzt?}{Czy był Pan już u lekarza?}
\turn{B}{Nein, ich konnte nicht kommen. Ich musste arbeiten.}{Nie, nie mogłem przyjść. Musiałem pracować.}
\turn{A}{Sie haben eine Erkältung. Nehmen Sie diese Tabletten und bleiben Sie zu Hause.}{Ma Pan przeziębienie. Proszę brać te tabletki i zostać w domu.}
\end{dialog}

\section{Ćwiczenia}
\begin{cwiczenia}
\exitem{Uzupełnij: \de{Ich \dots{} gestern krank.} (sein)}{\de{war}}
\exitem{Uzupełnij: \de{Du \dots{} Fieber.} (haben, Präteritum)}{\de{hattest}}
\exitem{Uzupełnij: \de{Wir \dots{} nicht kommen.} (können, Präteritum)}{\de{konnten}}
\exitem{Uzupełnij: \de{Er \dots{} zum Arzt gehen.} (müssen, Präteritum)}{\de{musste}}
\exitem{Uzupełnij: \de{Ihr \dots{} im Büro.} (sein, Präteritum)}{\de{wart}}
\exitem{Powiedz: „Boli mnie brzuch.” (z \de{wehtun})}{\de{Mir tut der Bauch weh.}}
\exitem{Powiedz: „Mam gorączkę.”}{\de{Ich habe Fieber.}}
\exitem{Powiedz: „Jestem chory od trzech dni.”}{\de{Ich bin seit drei Tagen krank.}}
\exitem{Zapytaj lekarza o dawkę: „Ile razy dziennie?” (zwrot: \de{Wie oft am Tag?})}{\de{Wie oft am Tag muss ich die Tabletten nehmen?}}
\exitem{Popraw: \de{Ich bin Schmerzen.}}{\de{Ich habe Schmerzen.}}
\exitem{Przetłumacz: „Wczoraj byłem chory i musiałem zostać w domu.”}{\de{Gestern war ich krank und musste zu Hause bleiben.}}
\exitem{Przetłumacz: „Zachorowałem.”}{\de{Ich bin krank geworden.}}
\end{cwiczenia}

\section{Trening słownictwa (generowany)}
\drill{17}

\weekend{Odegraj z AI wizytę lekarską: opisz objawy, odpowiedz na trzy pytania lekarza i powtórz zalecenia.
Wypisz 10 części ciała po niemiecku i przetestuj się na Anki (karty u17).}
